
To categorise the increasing level of complexity of DFAs in terms of the ingredients that entering the DFAs, as well as the modelling accuracies that are achievable with these DFAs, John Perdew and his colleague first introduced the `Jacob's ladder for DFAs'.\cite{perdew2001jacob} 

In the bottom three rungs of the Jacob's ladder, the general idea behind the development of DFAs is to increase its accuracy by gradually including more topological information on the single-particle electron density $\rho(\bm{r})$. With this, the exchange-correlation functional can be written in the general form of 
\begin{equation}\label{eq:Exc_general}
E_{xc}[\rho(\bm{r})]=\int (e_x^{F}[\rho(\bm{r})] + e_c^{F}[\rho(\bm{r})]) \mbox{d}\bm{r},
\end{equation}
where, starting from the exchange $e_{x}[\rho]$ and correlation $e_{c}[\rho]$ energies for homogeneous electron densities that can be either exactly determined or numerically parameterised, the exchange and correlation energies for inhomogeneous electron densities could be modelled by including an additional inhomogeneous factor $F$ that incorporates the information of electron density topologies, thus
\begin{equation}\label{eq:Exc_inhomo}
    e^{F}_{x,c}[\rho(\bm{r})]=F_{x,c}[\rho(\bm{r})] e_{x,c}[\rho],
\end{equation}
where the inhomogeneous factor can be written in a general form of:\cite{mardirossian2017thirty}
\begin{equation}\label{eq:inhomo_factor_general_form}
F_{IJ}[\rho(\bm{r})]=\sum_{i=0}^{I}\sum_{j=0}^{J} s_{ij}u_\sigma^i w_\sigma^j, 
\end{equation}
in which $s_{ij}$ are parameters. Dropping the explicit dependency on $\bm{r}$ for clarity, in \cref{eq:inhomo_factor_general_form}, the spin $(\sigma)$ dependent parameter $u_\sigma=\beta x_\sigma^2/(1+\beta x_\sigma^2)$ where $x_\sigma=|\nabla\rho_\sigma|/\rho_\sigma^{4/3}$ is so-called the reduced density gradient, and $w_\sigma=(t_\sigma-1)/(t_\sigma+1)$ where $t_\sigma=\frac{3}{5}(6\pi^2)^{\frac{2}{3}}\frac{\rho^{5/3}_\sigma}{\tau_\sigma}$ is related to the kinetic energy of a Kohn-Sham system $\tau_\sigma=\frac{1}{2}\sum_{i=1}^{\mbox{\scriptsize{occ}}}|\nabla \phi_{i,\sigma}^{\mbox{\scriptsize{KS}}}|^2$, which can be considered equivalent to measuring the curvatures of the electron density distributions. 

Using this formalism, the first rung of the Jacob's ladder contains exchange-correlation functionals that are constructed using $F_{00}$, which are functionals that only depend on the density of a homogeneous electron gas. This leads to the developments of \emph{Local Density Approximations} (LDA). The LDA has seen the initial success of DFAs in modelling condensed matters in which the electron densities vary slowly across the space. It provides good  values of the lattice constants albeit overestimates the binding energies of common inorganic solids.\cite{stampfl1999density} On the other hand, LDA did not win the confidence over DFAs in the chemistry community because its fundamental assumption of homogeneous electron density does not compile with the realistic behaviours of electrons involved in chemical bonds. This problem was overcome with the introduction of the \emph{Generalised Gradient Approximations} (GGA), by adding the reduced density gradients $x_\sigma$ into the recipe for constructing the DFAs that corresponds to the inclusion of $F_{10}$ term in \cref{eq:Exc_general}.  This has led to, for examples, the  PBE (Perdew-Burke-Ernzerhof)\cite{perdew1996generalized} functional that remains to be the workhorse in the condensed matter community as exemplified by the Materials Project,\cite{jain2013commentary} and the BP86\cite{becke1988density,perdew1986density}(augmented with dispersion correction) being a good choice in quantum chemistry to obtain good starting geometries for  more detailed analysis. Moving up in the Jacob's ladder,  with the $F_{11}$  term in \cref{eq:Exc_general}, one reaches the \emph{meta-GGA} functionals that include additionally, the electron kinetic energy in its construction recipe, which further improves the binding energies of molecular systems compared to pure-GGA.\cite{perdew1999accurate} 

As discussed, although direct combination of $E_{x}^{\mbox{\scriptsize{exact}}}$  with a correlation functional does not lead to good performing DFAs, the existence of a closed expression for $E_{x}^{\mbox{\scriptsize{exact}}}$ remains to be a highly attractive features for quantum chemists to explore different approaches to combine it with approximated correlation term to construct the \emph{hybrid functionals} that sit on the fourth rung of the Jacob's ladder. Notably this leads to the  hybrid functional B3LYP\cite{10.1063/1.464913,lee1988development,vosko1980accurate,stephens1994ab} that is most widely used in computational chemistry, although its stand-alone usage without  additional corrections are known to be not recommended for obtaining accurate results.\cite{kruse2012standard} Formally, the B3LYP functional is constructed\cite{stephens1994ab} as
\begin{equation}\label{eq:b3lyp_func}
    E_{xc}^{\mbox{\scriptsize{B3LYP}}}=(1-a_0)E_{x}^{\mbox{\scriptsize{LDA}}}+a_0 E_{x}^{\mbox{\scriptsize{exact}}}+a_x \Delta E_{x}^{\mbox{\scriptsize{B88}}}
    +a_c E_{c}^{\mbox{\scriptsize{LYP}}}+(1-a_c)E_{c}^{\mbox{\scriptsize{VWN}}}.
\end{equation}
Here, part of the LDA exchange is replaced by $E_{x}^{\mbox{\scriptsize{exact}}}$  with additional inclusion of Becke's gradient correction\cite{becke1988density} to the correlation functional. For the correlation energy, the Lee-Yang-Parr (LYP)\cite{lee1988development} functional, which lacks  an easily
separable local component, is combined with the local correlation
functional formulated by Vosko, Wilk, and Nusair (VWN).\cite{vosko1980accurate} $a_0$, $a_x$ and $a_c$ are three adjustable mixing parameters, which, in the case of B3LYP functional, does not contain explicit dependency on interatomic distance, such that functionals with this type of constructions are generally referred to as the \emph{global hybrid}. However, the Hartree-Fock exact exchange  describes the electronic interactions better than semilocal DFAs at large separations $r$ between electrons, whereas DFAs works better at short $r$'s due to Hartree-Fock's intrinsic lack of correlations between electrons with the same spins, this inspired the development of \emph{range-separated hybrids}, as pioneered by the CAM-B3LYP\cite{yanai2004new} functional that attenuates the mixing parameters with $r$. This  basic construct was modified by Heyd, Scuseria, and
Ernzerhof to circumvent the challenges of computing the exact exchange for periodic systems, leading to the HSE functionals\cite{heyd2003hybrid} as being the popular choice of hybrid functional in the condensed matter community.

Lastly, the \emph{double hybrid  functionals}, which are at the top rung of the Jacob's ladder, were proposed to overcome the inability for all lower-rung functionals in modelling the long-ranged dynamic electron correlation energies that play important contribution towards the non-covalent dispersion interactions in molecular system. Similar to how $E_{x}^{\mbox{\scriptsize{exact}}}$  is acquired in a hybrid DFA, the exchange-correlation energy from the double hybrid functional can be achieved by mixing part of the energy from hybrid DFA with the correlation energy obtained from a  second-order M\o ller-Plesset perturbation theory (MP2)\cite{moller1934note} calculation using both the occupied and virtual Kohn-Sham orbitals  $\left\{ \phi_{i}^{\mbox{\scriptsize{KS}}}\right\}$. Unlike DFAs, second-order perturbation theory belongs to the family of \emph{wavefunction theoretical methods} that are \emph{systematically improvable in terms of treating the challenging electron correlation problems in quantum chemistry,} which is beyond the scope of the present review. Using this strategy, \citet{grimme2006semiempirical} proposed the first double hybrid  functional dubbed as B2PLYP, in which the total exchange correlation energy is given by 
\begin{equation}\label{eq:B2PLYP_xc}
E_{xc}^{\mbox{\scriptsize{B2PLYP}}}=(1-a)E_{x}^{\mbox{\scriptsize{GGA}}}+a E_{x}^{\mbox{\scriptsize{exact}}}+bE_{c}^{\mbox{\scriptsize{GGA}}}+cE_{c}^{\mbox{\scriptsize{MP2}}},
\end{equation}
with $a$, $b$ and $c$ are the mixing parameters. Following this foundational constructions of double-hybrid exchange-correlation functionals, a range of different ideas across the chemistry community to construct similar functionals\cite{goerigk2014double} to achieve better accuracy in determining the energetics of molecular systems. Direct application of double hybrid DFAs in condensed matter simulations remains rare until these days due to the challenge in developing implementations to efficiently calculate the MP2 correlation energy with the planewave basis. \cite{carter2008challenges,pisani2008periodic}

\paragraph{\textbf{Parameterisation of DFAs}}
Although the goal of density functional \emph{theory} is to provide a formalism to solve the many-body Schr\"{o}dinger equation with no empirical parameters, which will enable us to predict the ground state energies and associated physical properties for new material systems, it can be seen above [\emph{e.g.} \cref{eq:b3lyp_func} and \cref{eq:B2PLYP_xc}] that some parameters are still necessary to be introduced in order to approximate the unknown exchange correlation functionals to make DFAs practical for solving real material problems. Typically, particularly in the early stages of DFA developments, these parameters are fitted to enable the total Kohn-Sham energies to reproduce the experimental thermochemical data. 

For examples, the popular B3LYP functional [\cref{eq:b3lyp_func}] was derived through a linear square fit to a dataset containing 56 atomisation energies, 42 ionisation potentials, 8 proton affinities and 10 first-row total atomic energies.\cite{beeke1993density} The mixing parameters in the B2PLYP functionals are fitted to the G2 dataset containing the heat of formation for 148 neutral molecules.\cite{curtiss1997assessment} The same dataset is also used for determining the Coulomb attenuation parameters in the CAM-B3LYP\cite{yanai2004new} functional. The hope here is to construct a DFA that can be universally applied to most chemical systems. However, this is also a monumental endeavour as the effects of electron correlations manifest themselves differently in different chemical systems (\emph{e.g.} main group molecules, transition metal complexes, non-covalent interactions, and so on), requiring system-specific functional forms and parameterisations of DFAs in order to achieve good performance. This is the motivation behind the development of the Minnesota family of density functionals,\cite{peverati2014quest,mardirossian2016accurate} and the associated database\cite{yu2016mn15} to optimise the parameters in the DFAs, depending on the specific (sub)sets of chemical problem(s) that a DFA is constructed to solve. 
